\begin{table}[H]
\centering
\caption{Descomposición del crecimiento del valor agregado bruto total, 2006--2024}
\label{tab:descomposicion_valor_agregado}
\small
\setlength{\tabcolsep}{12pt}
\begin{tabular}{lrrrr}
\toprule
 & Valor agregado bruto & Trabajo & Capital & PTF \\
 & (\% anual) & (pp por año) & (pp por año) & (\% anual) \\
\midrule
2006--2024 & 3,36 & 1,67 & 1,86 & -0,17 \\
\bottomrule
\end{tabular}
\vspace{0.2em}
\begin{minipage}{0.93\textwidth}
\footnotesize \textit{Nota:} el valor agregado bruto y la PTF se expresan como tasas de crecimiento logarítmicas anualizadas. Trabajo y capital son contribuciones al crecimiento, en puntos porcentuales por año, no tasas de crecimiento de los insumos. Como la PTF entra con coeficiente uno en la identidad contable, su tasa de crecimiento también equivale numéricamente a su contribución al crecimiento del valor agregado. Las cifras pueden no sumar por redondeo.
\par\textit{Fuente:} cálculos del CJC con base en DANE, anexo PTF 2025, Cuadro 1.
\end{minipage}
\end{table}
